\subsection{A finite and locally finite construction}

The regular-zero theorem is local and does not limit the number of connected
components globally.  The next model gives a direct answer to the existence part
of the multi-zero problem.

\begin{theorem}[Parallel multi-zero model]
\label{thm:parallel-multi-zero}
Let $h\in C^\infty(\R)$ have only simple zeros.  On $M=\R^2$ define
\begin{equation}
 a_h(x,y)=e^x h(y)
 \label{eq:parallel-assignment}
\end{equation}
and
\begin{equation}
 g_h=
 \frac{e^{2x}\bigl(h(y)^2+h'(y)^2\bigr)}
      {\mu^2+\lambda^2e^{2x}h(y)^2}
 (\dd x^2+\dd y^2).
 \label{eq:parallel-metric}
\end{equation}
Then $(M,g_h,a_h;\mu,\lambda)$ is a regular AES and
\begin{equation}
 Z(a_h)=\coprod_{h(c)=0}\R\times\{c\}.
 \label{eq:parallel-zero-set}
\end{equation}
In particular, if
$h_k(y)=\prod_{j=1}^k(y-c_j)$ with distinct $c_j$, the zero set has exactly
$k$ connected components.
\end{theorem}

\begin{proof}
One has
\[
 \dd a_h=e^x h(y)\,\dd x+e^x h'(y)\,\dd y,
 \qquad
 |\dd a_h|^2_{\mathrm{eucl}}
 =e^{2x}\bigl(h^2+(h')^2\bigr).
\]
The simple-zero hypothesis implies that $h$ and $h'$ never vanish
simultaneously.  The coefficient in \eqref{eq:parallel-metric} is therefore smooth
and positive.  Theorem~\ref{thm:conformal-realization} proves the AEG identity,
and $e^x$ never vanishes, so \eqref{eq:parallel-zero-set} is immediate.
\end{proof}

\begin{remark}[What the index counts]
\label{rem:descriptive-index}
The polynomial example proves existence of a regular AES with $k$ zero components.
It does not prove uniqueness, completeness, homogeneity, or a canonical
classification by $k$.  We therefore call it the \emph{parallel $k$-zero model},
not $E_k$.  The old repository symbol $E_k$ had no stable definition: at different
times its subscript suggested construction order, zero count, or singularity count.
\end{remark}

Taking $h(y)=\sin y$ gives a countable, locally finite zero family.  This example
already suggests that countably many visible sheets may be repetitions under a
covering group rather than independent components of downstairs geometry.

\subsection{A downstairs model and its logarithmic pullback}

Let $z=X+iY$ be the standard coordinate on $\C^\times$ and set
\begin{equation}
 a_\times(z)=Y,
 \qquad
 g_\times=\frac{\dd X^2+\dd Y^2}{\mu^2+\lambda^2Y^2}.
 \label{eq:punctured-plane-aes}
\end{equation}
Since $|\dd Y|^2_{\mathrm{eucl}}=1$, this is a regular AES.  Its zero set is the
punctured real axis, hence has two components.

Pull \eqref{eq:punctured-plane-aes} back along the exponential covering
\[
 \exp:\C\longrightarrow\C^\times,
 \qquad x+iy\longmapsto e^{x+iy}.
\]
The resulting data are
\begin{equation}
 \widetilde a(x,y)=e^x\sin y,
 \qquad
 \widetilde g=
 \frac{e^{2x}}
      {\mu^2+\lambda^2e^{2x}\sin^2y}
 (\dd x^2+\dd y^2),
 \label{eq:logarithmic-aes}
\end{equation}
and
\begin{equation}
 Z(\widetilde a)=\coprod_{k\in\Z}\{y=k\pi\}.
 \label{eq:logarithmic-zero-lines}
\end{equation}

\begin{theorem}[Logarithmic zero lift]
\label{thm:logarithmic-zero-lift}
The data in \eqref{eq:logarithmic-aes} form a regular AES and are invariant under
the deck transformation
\[
 D(x,y)=(x,y+2\pi).
\]
On the labels in \eqref{eq:logarithmic-zero-lines}, $D$ acts by $k\mapsto k+2$.
Thus the countable upstairs zero family has exactly two deck orbits, corresponding
to the positive and negative real half-axes downstairs.
\end{theorem}

\begin{proof}
The formula is the pullback of the regular AES
\eqref{eq:punctured-plane-aes}; it can also be checked directly from
$|\dd\widetilde a|^2_{\mathrm{eucl}}=e^{2x}$.  Both $\sin y$ and the metric
coefficient are $2\pi$-periodic.  Finally,
$D(\{y=k\pi\})=\{y=(k+2)\pi\}$, proving the orbit statement.
\end{proof}

\begin{remark}[No singular spine at the missing origin]
The coordinate $z=0$ is absent from $\C^\times$, but the formula
$a_\times=Y$ and its metric extend regularly across it.  Hence the origin is not a
locally essential singularity in the sense of Paper I.  The model is a
\emph{logarithmic-domain model}, not an essential-singularity model.  Calling the
missing point a singular spine would erase the local-essentiality condition.
\end{remark}

Theorem~\ref{thm:logarithmic-zero-lift} is the rigorous content we retain from the
historical $E_{\log}$ intuition:
\[
 \text{countably many sheets upstairs}
 \quad=\quad
 \text{finite downstairs geometry plus deck labels}.
\]
No claim is made that it reconstructs every meaning previously attached to
$E_{\log}$.

\subsection{Sheet labels are not automatically topological states}

The label $k\in\Z$ in \eqref{eq:logarithmic-zero-lines} depends on a choice of
fundamental domain.  Its parity is invariant under the deck group in this model,
but an absolute integer label is not downstairs data.  More generally, logarithms
of moving complex roots acquire integer shifts whose individual values depend on
the chosen lifts; only cycle sums survive the lift gauge.  This statement is proved
in Section~\ref{sec:braids} and Appendix~\ref{app:configuration}.

There is a second obstruction to interpreting the parallel model as a braid.  If
$c_1(t)<\cdots<c_k(t)$ are distinct real zero positions throughout a one-parameter
family, their order cannot change without a collision.  The configuration space of
$k$ unordered points on a line has contractible components; it does not carry the
Artin braid group.  Noncommutative braid monodromy requires motion in a surface,
which in this paper comes from complex-valued fields, or else a separately declared
transverse probe.

\subsection{Metric completeness and provenance limits}

The models above are exact solutions of the AEG equation, but completeness is not
part of their claim.  For example, in the parallel model the conformal factor often
decays like $e^{2x}$ as $x\to-\infty$, placing that end at finite metric distance.
Completion can add boundary or singular strata and must be audited separately.
This is one reason the model record must always list the domain and metric, not only
the zero set.

For reference, the construction record is:
\begin{center}
\small
\begin{tabular}{p{0.17\textwidth}p{0.34\textwidth}p{0.34\textwidth}}
\toprule
Item & Parallel model & Logarithmic model \\
\midrule
Domain & $\R^2$ & $\C\to\C^\times$ \\
Assignment & $e^xh(y)$ & $e^x\sin y=\im(e^{x+iy})$ \\
Metric & Equation~\eqref{eq:parallel-metric} &
Equation~\eqref{eq:logarithmic-aes} \\
Singular set & empty & empty \\
Zero topology & one line per simple zero of $h$ & countably many lines upstairs,
two components downstairs \\
Global status & regular; completeness not asserted & regular covering model;
missing origin removable \\
Provenance & constructed and proved here & constructed and proved here; historical
name used only as motivation \\
\bottomrule
\end{tabular}
\end{center}

\begin{problem}[Natural multi-zero metrics]
\label{prob:natural-multizero}
Classify multi-zero assignments when the metric is fixed by an independent AEG
construction, or when completeness, curvature, or expression-history naturality is
imposed.  Determine whether any canonical index survives those restrictions.
\end{problem}

Section~\ref{sec:arithmetic-zero-networks} gives one nonparallel answer without
solving this classification: an automorphic function fixed by a projective
arithmetic operator group determines both the metric and a four-valent infinite
zero network.  The contrast is intentional.  The present section proves that
multi-zero geometry is analytically abundant; the Hecke construction asks which
part of that abundance is forced by arithmetic structure.
